\chapter{Introduzione}
\section{Motivazione}
I dispositivi connessi intelligenti sono sempre più diffusi in abitazioni, laboratori e spazi industriali. La loro utilità aumenta quando riescono a interpretare i dati localmente; tuttavia, i vincoli di risorse e i requisiti di sicurezza rendono non banali le scelte progettuali \cite{shi2016edge,roman2018mobile}.

\begin{researchquestion}
Come può un dispositivo edge compatto offrire prestazioni di inferenza utili, mantenendo reattività e un confine di sicurezza chiaramente definito?
\end{researchquestion}

\section{Obiettivi}
Il progetto di laurea segue quattro obiettivi pratici:
\begin{itemize}
  \item definire un requisito di sistema misurabile;
  \item implementare un prototipo modulare;
  \item valutare latenza, accuratezza e uso delle risorse;
  \item discutere limiti e possibili miglioramenti.
\end{itemize}

\begin{keypoint}[Indicazioni per il modello]
Utilizzare il capitolo iniziale per presentare il contesto, definire il problema, dichiarare gli obiettivi e anticipare la struttura della tesi. Evitare dettagli implementativi che appartengono ai capitoli successivi.
\end{keypoint}

\section{Struttura del documento}
Il \cref{ch:b-design} presenta la progettazione, il \cref{ch:b-evaluation} descrive la valutazione e il capitolo finale riassume i risultati principali.

\chapter{Progettazione del sistema}\label{ch:b-design}
\section{Architettura}
Il dimostratore separa acquisizione, inferenza, politiche di sicurezza e presentazione. Il diagramma in \cref{fig:b-architecture} è generato direttamente in TikZ e può essere sostituito con una figura specifica del progetto.

\begin{figure}[htbp]
\centering
\begin{tikzpicture}[node distance=12mm and 9mm, every node/.style={font=\sffamily\small}]
\tikzset{block/.style={rounded corners=3pt,minimum width=31mm,minimum height=13mm,align=center,draw=ISISLine,fill=ISISPaper,text=ISISDeepBlue,font=\sffamily\bfseries}}
\node[block] (sensor) {Sensori\\e ingressi};
\node[block,right=of sensor] (edge) {Motore di inferenza\\edge};
\node[block,right=of edge] (policy) {Politica di\\sicurezza};
\node[block,below=of policy] (ui) {Dashboard\\e avvisi};
\draw[-{Latex[length=3mm]},line width=1.2pt,ISISLightBlue] (sensor) -- (edge);
\draw[-{Latex[length=3mm]},line width=1.2pt,ISISBlue] (edge) -- (policy);
\draw[-{Latex[length=3mm]},line width=1.2pt,ISISBlue] (policy) -- (ui);
\draw[-{Latex[length=3mm]},line width=1.2pt,ISISMuted,dashed] (ui.west) -| (edge.south);
\end{tikzpicture}
\caption{Architettura dimostrativa per la tesi di laurea.}
\label{fig:b-architecture}
\end{figure}

\section{Note implementative}
Descrivere hardware, moduli software, interfacce e principali scelte ingegneristiche. Per favorire la riproducibilità, registrare versioni e parametri di configurazione.

\begin{lstlisting}[language=Python,caption={Wrapper di inferenza dimostrativo.},label={lst:b-code}]
def classify(sample, model, threshold=0.75):
    score = model.predict(sample)
    return {
        "label": int(score >= threshold),
        "confidence": float(score),
    }
\end{lstlisting}

\begin{resultbox}[Scelta progettuale]
Il prototipo adotta una pipeline modulare, così da poter verificare in modo indipendente il modello di inferenza e la politica di sicurezza.
\end{resultbox}

\chapter{Valutazione sperimentale}\label{ch:b-evaluation}
\section{Protocollo}
Descrivere dataset o carico di lavoro, suddivisione train/test, baseline, metriche, hardware e prove ripetute. I valori riportati di seguito sono segnaposto.

\begin{figure}[htbp]
\centering
\begin{tikzpicture}
\begin{axis}[
width=0.84\textwidth,height=6.2cm,
ybar,bar width=17pt,
ymin=0,ymax=100,
ylabel={Punteggio (\%)},
symbolic x coords={Baseline,Prototipo,Ottimizzato},xtick=data,
axis line style={ISISLine},tick style={ISISLine},
ymajorgrids=true,grid style={ISISLine},
legend style={draw=none,fill=none,font=\sffamily\small},
label style={font=\sffamily\small,color=ISISMuted},
tick label style={font=\sffamily\small,color=ISISInk}]
\addplot[fill=ISISLightBlue,draw=ISISLightBlue] coordinates {(Baseline,72) (Prototipo,86) (Ottimizzato,91)};
\end{axis}
\end{tikzpicture}
\caption{Risultato di valutazione illustrativo. Sostituire con i dati del progetto.}
\label{fig:b-results}
\end{figure}

\begin{table}[htbp]
\centering
\caption{Esempio di sintesi delle prestazioni.}
\begin{tabularx}{0.82\textwidth}{Yrrr}
\toprule
Configurazione & Accuratezza & Latenza & Memoria\\
\midrule
Baseline & 72.0\% & 84 ms & 410 MB\\
Prototipo & 86.0\% & 41 ms & 286 MB\\
Ottimizzato & \textbf{91.0\%} & \textbf{29 ms} & \textbf{241 MB}\\
\bottomrule
\end{tabularx}
\label{tab:b-results}
\end{table}

\begin{keypoint}[Interpretazione]
Discutere dimensione dell'effetto, incertezza, rilevanza pratica, minacce alla validità e ogni discrepanza tra comportamento atteso e osservato.
\end{keypoint}

\chapter{Conclusioni}
Il progetto dimostra un percorso completo dalla definizione del problema all'implementazione e alla valutazione. Riassumere il contributo senza ripetere il sommario, quindi indicare limiti e sviluppi concreti.

\section{Sviluppi futuri}
Possibili estensioni includono un dataset più ampio, ulteriori modelli di minaccia, il deployment su hardware alternativo e uno studio utente dell'interfaccia di allerta.
